\section{Indexing Pipeline}
\label{sec:indexing-pipeline}

The indexing pipeline transforms a directory tree of source files into the
symbol graph stored in SQLite~\cite{brunsfeld2018treesitter}. It runs in five
stages organized into two phases: a structural phase that builds the symbol
graph and a deferred FTS phase that populates the content index.

\begin{figure}[htbp]
\centering
\begin{tikzpicture}[
  box/.style={draw, fill=green!8, rounded corners,
              minimum width=4.2cm, minimum height=0.9cm,
              text centered, font=\small\sffamily},
  phase/.style={draw=gray, dashed, rounded corners, inner sep=6pt},
  ann/.style={font=\scriptsize\ttfamily, text=gray},
  >=stealth, node distance=0.6cm
]
\node[box] (discover)  {File Discovery};
\node[box] (parse)     [below=of discover]  {tree-sitter Parse};
\node[box] (extract)   [below=of parse]     {Symbol \& Edge Extraction};
\node[box] (persist)   [below=of extract]   {Persist (batched SQLite)};
\node[box] (fts)       [below=0.9cm of persist]   {content\_fts Rebuild (phase 2)};

\draw[->] (discover) -- (parse);
\draw[->] (parse)    -- (extract);
\draw[->] (extract)  -- (persist);
\draw[->, dashed] (persist) -- (fts)
  node[ann, right, midway] {after structural commit};
\end{tikzpicture}
\caption{The codetopo indexing pipeline. FTS population is deferred to phase 2
         after the structural commit to avoid holding a write lock during
         content hashing.}
\label{fig:pipeline}
\end{figure}

\subsection{File discovery}

The pipeline begins by enumerating candidate files. When the workspace root
contains a \texttt{.git} directory, codetopo prefers \texttt{git ls-files} to
enumerate tracked files; this automatically respects \texttt{.gitignore} rules
and avoids scanning build artifacts. When no Git repository is detected, the
pipeline falls back to a recursive directory walk using the embedded
\texttt{.gitignore}-aware walker. Files are filtered by extension against the
registered set of language grammars.

Each candidate file is checked against its stored \texttt{mtime} and xxHash
fingerprint. If both match the stored values, the file is skipped entirely
(no parse, no write). This incremental check is the primary source of
startup-time savings on large repositories.

\subsection{Supported languages}

The following 13 languages are supported via tree-sitter grammars and
language-specific extractors: C, C++, C\#, TypeScript, JavaScript, Python,
Rust, Java, Go, Bash, SQL, YAML, and Lua.\footnote{Simon's architecture
summary lists 12; Lua is tracked as experimental.}

\subsection{Parsing}

Accepted files are parsed with the tree-sitter grammar for their detected
language (\texttt{src/index/parser.h:30--101}). Tree-sitter produces a
concrete syntax tree (CST) in memory. Parse errors are tolerated:
tree-sitter performs error-recovery and returns partial trees for files with
syntax errors. The CST is the input to the extractor.

\subsection{Symbol and edge extraction}

The extractor (\texttt{src/index/extractor.cpp:207--297}) walks the CST using
an \emph{iterative stack-based DFS} rather than recursion. This avoids stack
overflow on deeply nested ASTs (e.g., generated files with thousands of nested
expressions).

For each traversed node the extractor identifies:
\begin{itemize}
  \item \textbf{Symbol definitions}: functions, methods, classes, structs,
        enums, type aliases, variables, and constants. Each definition
        produces a \texttt{nodes} row with kind, name, qualified name,
        signature, source span, visibility, and optional docstring.
  \item \textbf{Call expressions}: recorded as \texttt{edges} of kind
        \texttt{calls}. Where the callee can be resolved to a known symbol
        in the current or previously indexed root, the edge carries both
        \texttt{src\_id} and \texttt{dst\_id}; unresolved calls carry
        \texttt{dst\_name} only and are flagged as approximate. A confidence
        floor of 0.3 is applied: edges below that threshold are discarded.
  \item \textbf{References}: import/include, inheritance, field access, and
        type references are recorded in the \texttt{refs} table.
\end{itemize}

\subsection{Persistence}

Extracted symbols, edges, and references are written to SQLite in
\emph{batched inserts} (\texttt{src/index/persister.h:117--380}):
\begin{itemize}
  \item Symbol (node) batch size: 100 rows per statement.
  \item Reference batch size: 80 rows per statement.
  \item Edge batch size: 150 rows per statement.
\end{itemize}

Before inserting, existing rows for the file are deleted
(\texttt{DELETE WHERE file\_id = ?}), so the operation is idempotent. The
entire insert sequence for one file runs inside a single transaction.

\paragraph{Bulk mode.} When the initial index covers more than 1\,000 files,
the pipeline enters \emph{bulk mode}: it drops secondary indexes and FTS
triggers before the load and rebuilds them afterwards. This avoids repeated
B-tree rebalancing during the insert flood and significantly reduces total
indexing time on large repositories.

\subsection{FTS rebuild (phase 2)}

After the structural commit (nodes, edges, refs) is finalized, the pipeline
runs a second phase that populates \texttt{content\_fts}. This table is a
\emph{contentless} FTS5 trigram index over source lines
(\S\ref{sec:data-model}). Because it is contentless, rows cannot be copied
from another database; they must be inserted by reading source files from
disk line-by-line. Deferring this phase means the write lock is held for only
as long as the structural data requires.
